\section{Statistics, limits and their common coordinates}\label{sec:setting}

\subsection{Model and statistics}
We observe $x_t=\mu+u_t$, $t=1,\dots,n$, where $(u_t)$ is stationary with mean zero and its partial sums satisfy the
functional central limit theorem $\sigma^{-1}n^{-H}\sum_{t\le\lfloor nr\rfloor}u_t\Rightarrow B_H(r)$ on $[0,1]$ for some
$0<H<1$ and $\sigma>0$. Here $B_H$ is fractional Brownian motion of Type~I \citep{MvN68}, the process with stationary
increments and covariance $\E B_H(s)B_H(t)=\frac12(s^{2H}+t^{2H}-|t-s|^{2H})$, so that $\Var B_H(1)=1$. We write
$Z=B_H(1)$ and $\Bt(r)=B_H(r)-rB_H(1)$ for the \emph{pathwise} bridge (not the conditioned bridge). Long memory
corresponds to $\frac12<H<1$, antipersistence to $0<H<\frac12$; $H=\delta+\frac12$ for the fractional integration
parameter $\delta$ of \citet{HI17} and $H=(\beta+1)/2$ for the memory parameter $\beta$ of \citet{MP12}.

Let $\bar x=n^{-1}\sum_tx_t$ and $S_k=\sum_{t\le k}(x_t-\bar x)$. The three statistics for $H_0:\mu=\mu_0$ are
\begin{align}
T_n&=\frac{n(\bar x-\mu_0)^2}{n^{-2}\sum_{k=1}^nS_k^2}, \label{eq:Tn}\\
\tau_{n,m}&=\frac{\sqrt n(\bar x-\mu_0)}{\hat\sigma_m},\quad
\hat\sigma_m^2=\frac{2\pi}m\sum_{j=1}^mI(\lambda_j),\quad I(\lambda)=\frac1{2\pi n}\Big|\sum_{t=1}^nx_te^{\mathrm it\lambda}\Big|^2,
\quad\lambda_j=\frac{2\pi j}n, \label{eq:tauHI}\\
t_{n,B}&=\frac{\sqrt n(\bar x-\mu_0)}{\hat\omega_B},\quad
\hat\omega_B^2=\frac1B\sum_{j=1}^B\hat\Lambda_j^2,\quad
\hat\Lambda_j=\sqrt{\frac2n}\sum_{t=1}^n\cos\Big(\frac{\pi j(t-\frac12)}n\Big)x_t. \label{eq:tEWC}
\end{align}
$T_n$ is the self-normalised statistic of \citet{Lob01} and \citet{Sha10}, and equals, up to the constant factors
recorded in \cref{tab:norm}, the Bartlett-kernel statistic with bandwidth equal to the sample size \citep{KVB00,KV02,MP12};
$\tau_{n,m}$ is equation~(2) of \citet{HI17}, without frequency weights; and $t_{n,B}$ is the EWC statistic, equation~(10)
of \citet{LLSW18}. Under $H_0$ and the conditions in these papers,
\begin{equation}\label{eq:limits}
\begin{gathered}
T_n\Rightarrow T_H=\frac{Z^2}{\int_0^1\Bt(r)^2\dd r},\\
\tau_{n,m}\Rightarrow\tau_{H,m}=\frac{Z}{\sqrt{m^{-1}\sum_{j=1}^m(C_j^2+S_j^2)}},\qquad
t_{n,B}\Rightarrow t_{H,B}=\frac{Z}{\sqrt{B^{-1}\sum_{j=1}^B\Lambda_j^2}},
\end{gathered}
\end{equation}
with the Wiener integrals
\[C_j=\int_0^1\cos(2\pi jr)\dd B_H(r),\qquad S_j=\int_0^1\sin(2\pi jr)\dd B_H(r),\qquad
\Lambda_k=\int_0^1\sqrt2\cos(k\pi r)\dd B_H(r).\]
At $H=\frac12$ the numerator is independent of the denominators and
$\tau_{1/2,m}\sim t_{2m}$, $t_{1/2,B}\sim t_B$; for $H\ne\frac12$ neither is true. A two-sided level-$\alpha$ test rejects
when $T_n>q_{1-\alpha}(H)$, $|\tau_{n,m}|>c_{m,\alpha}(H)$ or $|t_{n,B}|>c_{B,\alpha}(H)$, with $q_p(H)$ the $p$-quantile of
$T_H$ and $c_{\cdot,\alpha}(H)$ the upper $\alpha/2$-point of the symmetric limit.

\begin{table}[tbp]\centering\small
\caption{Normalisations of the published statistics and conversion of their tabulated values into quantiles $q_p$ of
$T_H$.}\label{tab:norm}
\footnotesize\setlength{\tabcolsep}{4pt}
\begin{tabular}{@{}>{\raggedright}p{0.21\textwidth}>{\raggedright}p{0.21\textwidth}>{\raggedright}p{0.31\textwidth}>{\raggedright\arraybackslash}p{0.19\textwidth}@{}}\toprule
source (location) & tabulated quantity & relation to $T_H$ & conversion\\\midrule
\citet{Lob01} (Table 1, $K=1$) & upper points of $U_1=T_n$ & $U_1\Rightarrow T_H$ & $q_p$ directly\\
\citet{AP97} (Table I, $k=1$) & quantiles $z_p$ of $Z_u$ & $Z_u\overset d=Z/(\int\Bt^2)^{1/2}$ \citep{AP02} & $q_{2p-1}=z_p^2$\\
\citet{KV02} (Table I) & quantiles $x_p$ of $t^\ast$ & $t^\ast=Z/(2\int\Bt^2)^{1/2}$ & $q_{2p-1}=2x_p^2$\\
\citet{MP09} (Tables 1--9, $b=1$) & one-sided points $x_a$ of $t_{\mathrm{MP}}$ & $t_{\mathrm{MP}}=Z/Q(1)^{1/2}$, $Q(1)=2\int\Bt^2$ & $q_{1-2a}=2x_a^2$\\
\citet{HI17} (Table 1) & two-sided $5\,\%$ points of $\tau$ & $\tau\Rightarrow\tau_{H,m}$ & $P(|\tau_{H,m}|>c)=0.05$\\
\citet{KPSS92} (Table 1) & upper points of $\eta_\mu$, $\eta_\tau$ & $\eta_\mu\Rightarrow\int_0^1\tilde W^2$ at $H=\frac12$ & bridge functional, validation\\
\bottomrule\end{tabular}
\end{table}

The normalisations matter in practice. The signed statistics of \citet{KV02} and \citet{MP12} satisfy $T_H=2t^2$, those of
\citet{KVB00} and \citet{AP97} satisfy $T_H=t^2$, and a one-sided test at level $a$ corresponds to the $(1-2a)$-quantile of
the squared statistic. Using a value from the wrong row of \cref{tab:norm} changes the level of the test by a factor
that is far larger than any error in the published tables; the audit of \cref{sec:audit} therefore matches each
published value to its own statistic, normalisation and tail before it converts it into a size.

\subsection{Common Gaussian coordinates}
All three limits in \eqref{eq:limits} are functionals of one Gaussian sequence. The following proposition, proved in
\cref{app:proofs}, makes this precise; it holds for every $0<H<1$.

\begin{proposition}[Cosine coordinates of the bridge]\label[proposition]{prop:coords}
Let $0<H<1$. Put $\Lambda_0=Z$ and, for $k\ge1$, $\Lambda_k=\int_0^1\sqrt2\cos(k\pi r)\dd B_H(r)$.
\begin{enumerate}
\item[(a)] For every $k\ge1$, $\int_0^1\Bt(r)\sqrt2\sin(k\pi r)\dd r=\Lambda_k/(k\pi)$ almost surely; consequently
\[\int_0^1\Bt(r)^2\dd r=\sum_{k\ge1}\frac{\Lambda_k^2}{(k\pi)^2}\quad\text{a.s. and in }L^1.\]
\item[(b)] The denominators in \eqref{eq:limits} are
$D_{\SN}=\sum_{k\ge1}\Lambda_k^2/(k\pi)^2$, $D_{\HI,m}=m^{-1}\sum_{j\le m}(\frac12\Lambda_{2j}^2+S_j^2)$ and
$D_{\EWC,B}=B^{-1}\sum_{j\le B}\Lambda_j^2$.
\item[(c)] $\Cov(\Lambda_k,\Lambda_l)=G_{kl}$ with the closed-form cosine Gram matrix of \citet[Proposition~1, closed-form Gram matrices]{AhmedCosChaos}:
with $b=2H-1$ and $p_k+\mathrm ir_k={}_1F_1(b;b+1;\mathrm ik\pi)$,
\begin{gather*}G_{00}=1,\qquad G_{0l}=-\frac{2\sqrt2Hr_l}{\pi l}\ (l\text{ even}),\qquad
G_{kk}=2H\Big(p_k-\frac{(1-b)r_k}{k\pi}\Big),\\
G_{kl}=\frac{4H}{\pi}\,\frac{lr_l-kr_k}{k^2-l^2}\ (k\ne l,\ k+l\text{ even}),\end{gather*}
and $G_{kl}=0$ when $k+l$ is odd.
\item[(d)] With $\Sigma_k=\int_0^1\sqrt2\sin(k\pi r)\dd B_H(r)$, so that $S_j=\Sigma_{2j}/\sqrt2$ and $C_j=\Lambda_{2j}/\sqrt2$,
\begin{gather*}\Cov(\Sigma_k,\Sigma_k)=2H\Big(p_k+\frac{(1+b)r_k}{k\pi}\Big),\\
\Cov(\Sigma_k,\Sigma_l)=\frac{4H}{\pi}\,\frac{kr_l-lr_k}{k^2-l^2}\ (k\ne l,\ k+l\text{ even}),\end{gather*}
$\Cov(\Sigma_k,\Sigma_l)=0$ when $k+l$ is odd, and $\Cov(\Sigma_i,\Lambda_k)=0$ whenever $i+k$ is even (in general not
otherwise). In particular every $S_j$ is uncorrelated with $Z$ and with every $C_l$.
\end{enumerate}
\end{proposition}

Part~(a) shows that the bridge sine coefficients are the EWC cosine coordinates divided by $k\pi$: the
Bartlett/self-normalised denominator is an EWC denominator with infinitely many terms and weights $(k\pi)^{-2}$, and the
fixed-$m$ denominator uses the even cosines together with the sines. The parity statements in~(c) follow from the
time-reversal symmetry $B_H(\cdot)\overset d=B_H(1)-B_H(1-\cdot)$, under which $\Lambda_k\mapsto(-1)^k\Lambda_k$,
$S_j\mapsto-S_j$ and $Z\mapsto Z$. In particular the fixed-$m$ vector $(Z,C_1,\dots,C_m,S_1,\dots,S_m)$ splits into the
two uncorrelated blocks $(Z,C)$ and $S$, and the odd cosines are uncorrelated with $(Z,\Lambda_2,\Lambda_4,\dots)$.

The formula in~(c) was derived in \citet{AhmedCosChaos} for $\frac12\le H<1$; it extends to antipersistence.

\begin{proposition}[Analytic continuation]\label[proposition]{prop:continuation}
The expressions for $G_{kl}$ in \cref{prop:coords}(c), and the corresponding sine Gram entries, are real-analytic in
$H\in(0,1)$ and equal $\Cov(\Lambda_k,\Lambda_l)$ for every $0<H<1$.
\end{proposition}

Consequently the fixed-$m$ and EWC limits are explicit finite-dimensional Gaussian ratio laws for every $0<H<1$. The
law of $T_H$ involves the infinite sum in~(a); its tail is controlled by the next theorem, a consequence of the two-sided
weighted cosine inequality of \citet[Lemma~3, two-sided weighted cosine inequality]{AhmedCosChaos}.

\begin{theorem}[Two-sided eigenvalue bounds for the fractional Brownian bridge]\label{thm:eig}
Let $\frac12\le H<1$ and $\cH=\Gamma(2H+1)\sin(\pi H)$. The eigenvalues $\lambda_1\ge\lambda_2\ge\dots$ of the covariance
operator $K_{\mathrm{br}}$ of $\Bt$ on $L^2(0,1)$ satisfy
\[\tfrac12\,\cH(k\pi)^{-1-2H}\le\lambda_k\le\cH(k\pi)^{-1-2H},\qquad k\ge1.\]
Moreover, for every $J\ge0$ the remainder $R_J=\sum_{k>J}\Lambda_k^2/(k\pi)^2$ is a Gaussian quadratic form whose
covariance eigenvalues are all at most $\bar c_J=\cH((J+1)\pi)^{-1-2H}$, and
$\E R_J=\tau_J:=\tr K_{\mathrm{br}}-\sum_{k\le J}G_{kk}/(k\pi)^2$ with $\tr K_{\mathrm{br}}=\frac1{(2H+1)(H+1)}-\frac16$.
\end{theorem}

At $H=\frac12$ the upper bound is attained, $\lambda_k=(k\pi)^{-2}$. Numerically $\lambda_k(k\pi)^{1+2H}/\cH\to1$ for
every $H$ we examined (\cref{fig:certificate}a), which suggests that the upper bound is asymptotically sharp, in line
with the exact asymptotics of \citet{CK18} for related fractional covariance operators.

\subsection{What is meant by size}
For a test with critical value $c$ we distinguish the \emph{exact finite-sample size}
$P_{\mu_0}(\text{reject at }n)$ under a fully specified Gaussian model, the \emph{limit size} under $T_H$, $\tau_{H,m}$ or
$t_{H,B}$, and the \emph{simulated} size of a Monte Carlo experiment. The critical value is the \emph{own limit} value
(the quantile of the limit at the true $H$), the \emph{naive} value (the quantile at $H=\frac12$: $q_p(\frac12)$, $t_{2m}$
or $t_B$), or a \emph{plug-in} value at an estimate $\hat H$. Following the conventions frozen before any computation
(\cref{app:audit}), a value is \emph{exact} when it is a closed form or a finite-dimensional law evaluated with a stated
numerical error, \emph{certified} when every numerical step is enclosed with outward rounding, and \emph{computed} when
it is a floating-point result without a complete enclosure.
